\scanpage{3}ment de Coxeter\inserted{\handwritten{$c$}} de $W$ qui la normalise. Alors le normalisateur de $\deleted{W}D$ dans $W$ est le groupe engendré par $c$, et il existe une unique élément de Coxeter de $W$ normalisant $D$ et y ayant pour valeur propre une racine \deleted{de}primitive \unsure{\inserted{$h$.ème}} donnée de $1$. Les droites régulières $D$ ayant la propriété précédente (appelées ``droites cycliques'' par Kostant) sont conjuguées sous $W$, leur nombre est égal à $\operatorname{card}(W)/h$.

NB Dans l'ensemble des droites cycliques, le groupe $(\underline{Z}/h\underline{Z})^{\handwritten{*}}$ opère\inserted{en commutant à l'action de $W$} librement; deux droites cycliques sont conjuguées sous l'action de ce groupe si et seulement si elles ont même stabilisateur dans $W$.
% archive sheet number: 2
